%%%%%%%%%%%%%%%%%%%%%%%%%%%%%%%%%%%%%%%%%%%%%%%%%%%%%%%%%%%%%%%%%%%%%%%%%%%%%%%
%% Guia de Aulas Práticas -- documento principal
%%
%% Compilar com:  latexmk -pdf main.tex
%% (duas passadas são necessárias por causa do total de páginas no rodapé)
%%%%%%%%%%%%%%%%%%%%%%%%%%%%%%%%%%%%%%%%%%%%%%%%%%%%%%%%%%%%%%%%%%%%%%%%%%%%%%%
\documentclass{guiapratico}

%  ===== Metadados do documento
%% Dados do guia -- edite apenas este arquivo para adaptar a outra disciplina.

\universidade{UNIVERSIDADE FEDERAL DE SÃO JOÃO DEL-REI}
\subtitulos{%
	CAMPUS ALTO PARAOPEBA \\
	PRÓ-REITORIA DE ENSINO DE GRADUAÇÃO -- PROEN\\
	COORDENADORIA DO CURSO DE ENG. MECATRÔNICA -- CEMEC}

\titulo{GUIA DE AULAS PRÁTICAS}

\unidadecurricular{Laboratório de Controle de Sistemas}
\cursos{Eng. Mecatrônica}
\professor{Edgar Campos Furtado, Filipe Augusto Santos Rocha}

\arquivo{}
\datadoc{}
\revisao{}

\localano{Ouro Branco, 2026}

%  ===== Bibliografia
%% Base de referencias no padrao ABNT (NBR 6023:2018). Pode ser chamado mais de
%% uma vez, se houver mais de um arquivo .bib.
\bibliografia{referencias.bib}



% ===== Documento
\begin{document}

\maketitle

% --- sumário (chapter > section > subsection > subsubsection)
\sumario

% --- Introdução
\secaoguia{Introdução}

O laboratório de \nomeunidade\ tem por objetivo complementar e sedimentar o
aprendizado dos conceitos teóricos apresentados em sala de aula. Para tanto, estão
previstas as aulas práticas listadas a seguir:

\begin{enumerate}
	\item Portas Lógicas;
	\item Teoremas Booleanos --- Parte 1;
	\item Circuitos Combinacionais;
\end{enumerate}

No Apêndice~A são apresentadas as normas para uso do laboratório, as quais visam
organizar o uso do mesmo, bem como prover uma maior segurança aos alunos durante a
realização das aulas práticas. Pede-se que seja feita uma leitura cuidadosa dessas
normas.

O material teórico de apoio às práticas é o adotado na disciplina teórica
correspondente \cite{ogata2010,nise2017}. Para os fundamentos de eletrônica digital
empregados nas primeiras aulas, recomenda-se \textcite{tocci2018}; as referências dos
relatórios devem seguir a norma \cite{nbr6023}.

As aulas práticas foram planejadas para serem realizadas no tempo de \qty{1}{\hour} e \qty{40}{\minute}, que
corresponde a duas horas/aula. Exige-se um relatório sobre a prática realizada, que
deve ser entregue na aula subsequente. A organização do relatório, bem como a capa do
mesmo, é apresentada no Apêndice~B.

% --- aulas
%%%%%%%%%%%%%%%%%%%%%%%%%%%%%%%%%%%%%%%%%%%%%%%%%%%%%%%%%%%%%%%%%%%%%%%%%%%%%%%
%% Aula 1 -- Ambientação com a planta TQ CE117 e conexão via pycomm3
%%
%% Local: aulas/aula-01/conteudo.tex
%% Imagens: aulas/aula-01/img/{planta.jpg,diagrama-processo.png}
%% Códigos: aulas/aula-01/src/{conversoes.py,ler_planta.py,comanda_planta.py}
%%
%% Material de origem: página do Trabalho Final da planta TQ CE117 (Kuara/UFSJ),
%% de onde vêm o diagrama de processo, as tabelas de instrumentação, os dados
%% do CLP e os parâmetros de rede. Os scripts de referência estão em
%% https://github.com/profilgusto/clp-abs-comms
%%%%%%%%%%%%%%%%%%%%%%%%%%%%%%%%%%%%%%%%%%%%%%%%%%%%%%%%%%%%%%%%%%%%%%%%%%%%%%%
\aula[Ambientação com a planta TQ CE117]{Ambientação com a Planta TQ CE117 e Comunicação com o CLP em Python}

\objetivos{%
  \item identificar no diagrama de processo os dois circuitos hidráulicos da planta TQ CE117 e a função de
    cada instrumento;
  \item relacionar cada instrumento à \sw{tag} correspondente no CLP e à sua faixa de conversão A/D ou D/A;
  \item estabelecer, em Python, a comunicação \sw{EtherNet/IP} com o CLP da planta usando a biblioteca
    \sw{pycomm3};
  \item converter as contas lidas do CLP em grandezas físicas e as porcentagens de acionamento em contas de
    saída;
  \item operar o circuito frio aberto por linha de comando e por interface gráfica, enchendo o tanque de
    processo com segurança;
  \item levantar experimentalmente a curva de calibração do transmissor de nível, relacionando a tensão lida
    à altura de líquido medida com régua milimetrada;
  \item reunir as rotinas de leitura, escrita e conversão em um módulo Python reaproveitável, base dos
    algoritmos de controle das aulas seguintes.}

\listamaterial{1 planta TQ CE117 (Laboratório de Controle e Automação, sala 103.4); 1 CLP Allen-Bradley com
  CPU CompactLogix L24EPR, já programado e energizado; 1 computador com Python 3.8 ou superior e acesso ao
  roteador \sw{Wi-Fi} do laboratório; 1 régua milimetrada (ou trena com escala em milímetros) para medir a altura de líquido no tanque de processo;
  água limpa no reservatório da planta.}

\campo{Duração:}{2 horas/aula (\qty{1}{\hour} e \qty{40}{\minute}).}

\section{Introdução}

Um sistema de controle normalmente se organiza em camadas.
%
Na base está o processo físico equipado com diversos \textbf{instrumentos}, classificados em dois grupos distintos:
%
\begin{itemize}
  \item os \textbf{sensores} permitem converter grandezas físicas em sinais elétricos;
  \item os \textbf{atuadores} são capazes de aplicar energia, agindo assim sobre o processo.
\end{itemize}
%
Acima dos instrumentos está o controlador, responsável pela aquisição determinística dos sinais e pelo acionamentos dos atuadores.
%
Estando estes três elementos interligados via rede de comunicação, torna-se possível implementar um controle retro-alimentado no processo, ou seja, realizar ações na planta em função dos dados provenientes dos sensores a partir de uma lógica pré-programada.


Neste contexto, o objetivo desta aula é introduzir a planta didática que utilizaremos nesta disciplina, descrevendo seu funcionamento e instrumentos, além do controlador utilizado e maneiras de comunicar e interagir de maneira programática com estes elementos. 
%
Ao final desta prática, será possível transformar as leituras dos sensores localizados na planta em variáveis de memória em um algoritmo python; e, vice-versa, será possível programar comandos para os atuadores neste algoritmo que serão convertidos em ações dos atuadores na planta.
%
Esta prática é a base de todo o resto que será construído em termos de modelagem e controle da planta didática, sendo então o domínio da prática atual essencial para o bom andamento do curso.


\subsection{A planta TQ CE117}
Iremos utilizar a planta didática TQ CE117 (\rfig{fig:planta}), que emula um processo industrial contínuo \cite{tecquipment2008}.
%
A planta possui um painel de processo que interagem diretamente com os instrumentos da planta e que disponibiliza sinais analógicos em seus contatores elétricos.
%
Um CLP Allen-Bradley está conectado a esta interface através de cartões de conversão A/D e D/A.


O elemento central desta planta é o tanque de processo, um cilindro transparente com escala graduada, através do qual se bombeia água por dois circuitos hidráulicos distintos (apresentados no diagrama de processo da \rfig{fig:diagrama}):
\begin{itemize}
  \item \textbf{Circuito quente fechado:} \cd{PUMP1} recircula água entre o tanque aquecedor --- onde está o aquecedor elétrico --- e o trocador de calor imerso no tanque de processo. A vazão é medida por \cd{FT1} e a temperatura do tanque aquecedor, por \cd{TT1}. Nenhuma água desse circuito se mistura à do tanque de processo: a troca é apenas térmica.
  \item \textbf{Circuito frio aberto:} \cd{PUMP2} recalca água do reservatório, passando pelo radiador (com ventoinha e trocador de calor) e pela válvula proporcional \cd{S}, até a entrada do tanque de processo. A vazão é medida por \cd{FT2}, e as temperaturas de entrada e saída do radiador, por \cd{TT3} e \cd{TT4}. A
    água retorna ao reservatório pela válvula de dreno, na base do tanque.
\end{itemize}

O tanque de processo é instrumentado com o transmissor de pressão \cd{PT}, o transmissor de nível \cd{LT} e o transmissor de temperatura \cd{TT5}. Uma válvula de ventilação, no topo, e uma válvula \sw{by-pass}, no retorno do circuito frio, são operadas manualmente.

\begin{figure}[H]
  \centering
  \includefigure[0.62\textwidth]{planta.jpg}
  \caption{Planta TQ CE117 do Laboratório de Controle e Automação.}
  \label{fig:planta}
\end{figure}

\begin{figure}[H]
  \centering
  \includefigure[0.85\textwidth]{diagrama-processo.png}
  \caption{Diagrama de processo da planta TQ CE117.}
  \label{fig:diagrama}
\end{figure}

Neste curso porém nos concentraremos somente no tanque de processo e no circuito frio aberto, buscando implementar leis de controle que regulam a altura do nível da água neste tangue.


\subsection{Instrumentação}
Todos os sensores e atuadores da planta utilizam sinal em tensão análogica de \qtyrange{0}{10}{\volt}, sendo seus contatos localizados no painel de processo localizado ao lado da planta didática.
%
As Tabelas~\ref{tab:entradas} e \ref{tab:saidas} apresentam as conversões \textbf{nominais} de cada instrumento do ponto de vista do CLP.

\begin{table}[H]
  \centering
  \caption{Entradas: sensores da planta e conversão para unidade de engenharia.}
  \label{tab:entradas}
  \begin{tabelaguia}{colspec={llcX[l]},width=\textwidth}
    Instrumento & \sw{Tag} & Sinal & Conversão \\
    Sensores de vazão & FT1, FT2 & \qtyrange{0}{10}{\volt} & \qty{1}{\litre\per\minute} por volt \\
    Sensores de temperatura & TT1 a TT5 & \qtyrange{0}{10}{\volt} & \qty{10}{\celsius} por volt \\
    Sensor de nível & LT & \qtyrange{0}{10}{\volt} & \qty{0}{\volt}: tanque vazio; \qty{10}{\volt}: tanque cheio \\
    Sensor de pressão & PT & \qtyrange{0}{10}{\volt} & \qty{100}{\milli\bar} por volt \\
  \end{tabelaguia}
\end{table}

\begin{table}[H]
  \centering
  \caption{Saídas: atuadores da planta e conversão do comando.}
  \label{tab:saidas}
  \begin{tabelaguia}{colspec={llcX[l]},width=\textwidth}
    Instrumento & \sw{Tag} & Sinal & Conversão \\
    Aquecedor elétrico & --- & \qtyrange{0}{10}{\volt} & \qty{75}{\watt} por volt \\
    Válvula proporcional & S & \qtyrange{0}{10}{\volt} & \qty{0}{\volt}: fechada; \qty{10}{\volt}: totalmente aberta \\
    Bombas de água & PUMP1, PUMP2 & \qtyrange{0}{10}{\volt} & \qty{0}{\volt}: sem fluxo; \qty{10}{\volt}: fluxo máximo \\
  \end{tabelaguia}
\end{table}

\subsection{O CLP e a conversão analógica-digital}
Os instrumentos estão ligados aos cartões analógicos de um CLP Allen-Bradley com CPU CompactLogix L24EPR \cite{rockwell-compactlogix}.
%
Esses cartões trabalham na faixa de \qtyrange{-10,5}{+10,5}{\volt}, e a conversão A/D e D/A é representada na memória do CLP por um inteiro (\cd{INT}) de \qty{16}{\bits}, sendo um de sinal (i.e., positivo/negativo) e quinze de valor.
%
Esse inteiro é comumente chamado de \textbf{conta}: \qty{-10,5}{\volt} corresponde a \num{-32768} contas e \qty{+10,5}{\volt}, a \num{+32767} contas.
%
A tensão do instrumento pode ser calculada a partir da conta $N$ lida pelo cartão do CLP através da equação

\begin{equation}
  V = \frac{10{,}5}{32768}\,N \qquad [\unit{\volt}],
  \label{eq:adc}
\end{equation}

\noindent e a conta a ser escrita para produzir uma tensão $V$ de comando é

\begin{equation}
  N = \operatorname{round}\!\left(\frac{32768}{10{,}5}\,V \right),
  \label{eq:dac}
\end{equation}

onde $\operatorname{round(\cdot)}$ representa uma função programática que arredonda um número float para o inteiro mais próximo.

\obs{Os instrumentos trabalham de \qtyrange{0}{+10}{\volt}, enquanto os cartões cobrem
\qtyrange{-10,5}{+10,5}{\volt}; logo, o fundo de escala do cartão não é o fundo de escala do instrumento.
%
Por exemplo, um sensor saturado não indica \num{32767} em sua leitura, mas sim $~$\num{31207}.
%
Usar o número errado nessa conta é o erro mais comum na primeira conexão com a planta.}

\obs{o maior valor que cabe num \texttt{INT} é \num{32767}. Escrever \num{32768} numa \sw{tag} de saída faz o
valor transbordar e chegar ao CLP como \num{-32768}, ou seja, comando de \qty{-10,5}{\volt} em vez de fundo de
escala. Toda escrita deve ser saturada em \num{32767} antes de sair do programa em Python.}

O programa já carregado na memória no CLP expõe à rede o subconjunto de \textit{tags} da Tabela~\ref{tab:tags}, sendo somente estas acessíveis pelo seu programa.

\begin{table}[H]
  \centering
  \caption{\sw{Tags} do CLP acessíveis pela rede.}
  \label{tab:tags}
  \begin{tabelaguia}{colspec={lccX[l]},width=\textwidth}
    \sw{Tag} & Tipo & Faixa & Função \\
    Program:MainProgram.PUMP2\_DAC  & INT & \numrange{0}{32767} & comando da bomba do circuito frio \\
    Program:MainProgram.VALVE\_DAC  & INT & \numrange{0}{32767} & comando da válvula proporcional S \\
    Program:MainProgram.FT2\_ADC    & INT & \numrange{-32768}{32767} & vazão no circuito frio aberto \\
    Program:MainProgram.PT\_ADC     & INT & \numrange{-32768}{32767} & pressão no tanque de processo \\
    Program:MainProgram.TT5\_ADC    & INT & \numrange{-32768}{32767} & temperatura no tanque de processo \\
    Program:MainProgram.LT\_ADC     & INT & \numrange{-32768}{32767} & nível no tanque de processo \\
  \end{tabelaguia}
\end{table}

Neste curso, iremos focar apenas nos instrumentos ligados ao monitoramento e controle do nível de vazão de água no tanque, mais especificamente (\textit{tags} abreviadas): \cd{PUMP2_DAC}, \cd{VALVE_DAC}, \cd{FT2_ADC} e \cd{LT_ADC}.


\subsection{Interagindo com a planta usando Python}
A comunicação com o CLP será feita diretamente em Python, pela biblioteca \sw{pycomm3}, que implementa o protocolo \sw{EtherNet/IP} sobre o protocolo CIP e conversa com controladores Logix sem exigir \sw{software} proprietário \cite{pycomm3,odva-cip}. 

Toda a interação do código Python com o CLP pode ser resumida a três chamadas básicas:
%
\begin{itemize}
  \item \cd{LogixDriver(ip)} --- abre a sessão CIP com o controlador; usado como gerenciador de contexto
    (\cd{with}), fecha a sessão sozinho ao final do bloco;
  \item \cd|plc.read('tag1', 'tag2', ...)| --- lê uma ou mais \sw{tags} de uma vez e devolve uma lista de
    resultados, cada um com os campos \cd{tag}, \cd{value}, \cd{type} e \cd{error};
  \item \cd|plc.write(('tag', valor), ...)| --- escreve nas \sw{tags} indicadas e devolve resultados no mesmo
    formato, para que a confirmação do CLP seja verificada.
\end{itemize}

O Código~\ref{cod:leitura-minima} ilustra um programa mínimo que utiliza esta interface, abrindo a sessão, lendo uma \textit{tag} em seu valor \cd{INT} e printando o resultado.

\begin{codigopython}[caption={Programa mínimo de leitura de uma \texttt{tag} com o \texttt{pycomm3}.},
                     label={cod:leitura-minima}]
from pycomm3 import LogixDriver

with LogixDriver('200.200.200.25') as plc:
    r = plc.read('Program:MainProgram.LT_ADC')
    print(r.tag, '=', r.value, 'contas')
\end{codigopython}


Um conjunto de exemplos prontos para esta planta está disponível no repositório Github \cd{clp-abs-comms}\footnote{\url{https://github.com/profilgusto/clp-abs-comms}}, organizado em três níveis:

\begin{itemize}
  \item \cd{1_interacao-basica/} --- \cd{ler_tags.py}, que lê todas as \sw{tags} da planta, e
    \cd{atua_e_le.py}, que escreve nos DACs, aguarda \qty{2}{\second} e lê o resultado de volta;
  \item \cd{2_gui/} --- \cd{gui_planta.py}, interface gráfica em \sw{tkinter} com botões liga/desliga para
    \cd{PUMP2} e \cd{VALVE} e gráfico do nível em tempo real; a opção \cd{--sim} simula o tanque, para testar
    o programa longe do laboratório;
  \item \cd{3_interacao-matlab/} --- \cd{daemon_clp.py}, que mantém a sessão CIP aberta e a expõe numa porta
    TCP em JSON, permitindo que o \sw{MATLAB} leia e escreva as mesmas \sw{tags} sem nenhuma \sw{toolbox}.
\end{itemize}

Os códigos desta aula, listados ao final da prática, seguem a mesma organização: \cd{conversoes.py} concentra
as conversões das \req{eq:adc} e \req{eq:dac}, e é importado por \cd{ler_planta.py} (aquisição) e
\cd{comanda_planta.py} (acionamento).

Os dados de rede do laboratório são: SSID \cd{Lab_Automa}, senha \cd{devibus485} e CLP no endereço
\cd{200.200.200.25}.

\obs{Segurança da planta: a válvula proporcional \texttt{S} deve estar totalmente aberta \emph{antes} de
qualquer acionamento de \texttt{PUMP2}. Acionar a bomba contra a válvula fechada pressuriza a linha e pode
danificar o equipamento. Esse intertravamento é responsabilidade do programa em Python que você vai escrever
--- o CLP não o implementa.}

\section{Prática}

\begin{roteiro}
  \item Percorra a planta e localize fisicamente, comparando com a \rfig{fig:diagrama}:
    \begin{roteiro}
      \item o tanque de processo, o reservatório e o tanque aquecedor;
      \item \cd{PUMP1}, \cd{PUMP2}, o radiador com a ventoinha e o trocador de calor;
      \item a válvula proporcional \cd{S}, a válvula \sw{by-pass}, a válvula de ventilação e a válvula de
        dreno inferior;
      \item os transmissores \cd{LT}, \cd{PT} e \cd{TT5}, junto ao tanque de processo, e \cd{FT2}, na linha do
        circuito frio.
    \end{roteiro}
  \item Verifique que a água está acima do nível mínimo no reservatório inferior.
  \item Abra a válvula de dreno inferior do tanque de processo e aguarde o tanque esvaziar completamente para
    o reservatório. Mantenha a válvula de ventilação aberta.
  \item Conecte o computador à rede \cd{Lab_Automa} e confirme o alcance do CLP com um \cd{ping} para
    \cd{200.200.200.25}. Registre o tempo de resposta médio.
    \obs{Sem resposta ao \texttt{ping} não há o que depurar no programa: o problema é de rede, não de código.}
  \item Prepare o ambiente Python, clonando o repositório de exemplos e instalando a biblioteca em um ambiente
    virtual:
    \begin{roteiro}
      \item \cd|git clone https://github.com/profilgusto/clp-abs-comms|;
      \item \cd|python3 -m venv .venv| e \cd|source .venv/bin/activate|;
      \item \cd|pip install -r 1_interacao-basica/requirements.txt|.
    \end{roteiro}
  \item Execute \cd|python3 1_interacao-basica/ler_tags.py| e confirme que o nome do controlador é
    impresso e que as seis \sw{tags} da Tabela~\ref{tab:tags} respondem sem erro. Registre na
    Tabela~\ref{tab:leituras} o valor em contas de cada sensor, com o tanque ainda vazio.
    \obs{o \texttt{pycomm3} costuma imprimir alguns avisos de depreciação antes da saída útil. Eles são
    inofensivos; para suprimi-los, execute o script com a opção \texttt{-W ignore} do interpretador.}
  \item Abra o Código~\ref{cod:conversoes} e confira, linha a linha, como as \req{eq:adc} e \req{eq:dac} foram
    implementadas e onde entra o ganho de cada instrumento da Tabela~\ref{tab:entradas}.
  \item Execute \cd|python3 ler_planta.py| (Código~\ref{cod:leitura}) e complete a
    Tabela~\ref{tab:leituras} com a tensão e o valor em unidade de engenharia exibidos. Calcule à mão, pela
    \req{eq:adc}, a tensão correspondente a uma das contas registradas e compare-a com a exibida pelo
    programa.
  \item Abra a válvula proporcional \cd{S} por inteiro com
    \cd|python3 comanda_planta.py 100 0| (Código~\ref{cod:comando}) e confirme, pelo ruído do
    atuador e pela inspeção visual, que ela abriu totalmente.
  \item Verifique o intertravamento antes de usá-lo de verdade: execute
    \cd|python3 comanda_planta.py 0 25| e confirme que o programa recusa o comando e não escreve
    nada no CLP.
  \item Com a válvula \cd{S} aberta em \qty{100}{\percent}, feche a válvula de dreno inferior e acione a bomba
    em \qty{25}{\percent} com \cd|python3 comanda_planta.py 100 25|.
  \item Aguarde a estabilização da vazão e, com \cd|python3 ler_planta.py 10 0.5|, registre na
    Tabela~\ref{tab:enchimento} o comando aplicado, a vazão \cd{FT2} indicada e a leitura de \cd{LT}. Anote
    também a altura de líquido medida com a régua milimetrada.
  \item Repita os dois passos anteriores para comandos de \qty{50}{\percent} e \qty{75}{\percent} em
    \cd{PUMP2}, sem deixar o nível ultrapassar \qty{80}{\percent} da escala do tanque.
    \obs{Acompanhe o tanque visualmente durante todo o enchimento. O transmissor de nível não aciona
    intertravamento algum: o único limite é o operador.}
  \item Prepare a calibração do transmissor de nível: com a válvula de dreno ainda fechada, encha o tanque
    até aproximadamente \qty{80}{\percent} da sua altura e então pare a bomba com
    \cd|python3 comanda_planta.py 100 0|. Aguarde a superfície do líquido ficar parada.
    \obs{a leitura de \texttt{LT} só é confiável com o líquido em repouso: com a bomba ligada, a agitação na
    entrada do tanque e a coluna em movimento deslocam a indicação. Meça sempre com vazão nula.}
  \item Registre o primeiro ponto da Tabela~\ref{tab:calibracao-lt}: encoste a régua na parede do tanque, com
    o zero no fundo interno, e leia a altura $h$ da superfície do líquido, em milímetros; em seguida execute
    \cd|python3 ler_planta.py| e anote as contas e a tensão de \cd{LT}.
    \obs{leia a régua com o olho na altura da superfície, para evitar erro de paralaxe, e sempre no mesmo
    ponto do tanque. O menisco deve ser lido pela sua base, como em uma proveta.}
  \item Levante os demais pontos da Tabela~\ref{tab:calibracao-lt} esvaziando o tanque aos poucos:
    \begin{roteiro}
      \item abra a válvula de dreno inferior o suficiente para um escoamento lento e feche-a novamente após
        o nível cair cerca de \qtyrange{5}{10}{\milli\metre} --- o passo exato não importa, desde que os
        pontos fiquem distribuídos por toda a faixa;
      \item espere a superfície estabilizar, meça $h$ com a régua e registre a leitura de \cd{LT};
      \item repita até o tanque esvaziar, completando as \num{20} linhas da tabela, e inclua como último
        ponto o tanque vazio ($h = \qty{0}{\milli\metre}$).
    \end{roteiro}
    \obs{a válvula de dreno deve estar fechada no instante da medição. Medir com o tanque escoando introduz
    erro nos dois sentidos: a altura muda enquanto se lê a régua e a leitura de \texttt{LT} acompanha com
    atraso.}
  \item Confira, ainda na bancada, se os pontos coletados são coerentes: a tensão deve crescer de forma
    monotônica com a altura. Se algum ponto destoar visivelmente dos vizinhos, repita a medição antes de
    desmontar o ensaio.
  \item Feche a válvula \cd{S} com \cd|python3 comanda_planta.py 0 0| e confirme que ambas as saídas ficaram
    zeradas.
  \item Execute a interface gráfica do repositório com \cd|python3 2_gui/gui_planta.py| e repita um
    enchimento curto usando os botões liga/desliga, acompanhando a curva de \cd{LT} na tela. Compare a leitura
    mostrada no gráfico com a impressa pelo seu próprio programa.
    \obs{a GUI zera as saídas ao ser fechada, desde que a opção \emph{zerar saídas ao sair} esteja marcada.
    Confirme visualmente que a bomba parou antes de deixar a bancada.}
  \item Guarde os três arquivos Python da aula (\cd{conversoes.py}, \cd{ler_planta.py} e
    \cd{comanda_planta.py}) --- eles serão o ponto de partida das aulas seguintes, nas quais o laço de
    controle será fechado sobre as mesmas funções de leitura e escrita.
\end{roteiro}


\subsection{Tabelas}

\begin{table}[H]
  \centering
  \caption{Leituras de referência com o tanque vazio --- preencher durante a prática.}
  \label{tab:leituras}
  \begin{tabelaguia}{colspec={cccc},width=0.9\textwidth}
    Instrumento & Contas lidas & Tensão [\unit{\volt}] & Unidade de engenharia \\
    FT2 & \vazio & & \\
    LT  & \vazio & & \\
    PT  & \vazio & & \\
    TT5 & \vazio & & \\
  \end{tabelaguia}
\end{table}

\begin{table}[H]
  \centering
  \caption{Enchimento do tanque pelo circuito frio aberto --- preencher durante a prática.}
  \label{tab:enchimento}
  \begin{tabelaguia}{colspec={ccccc},width=\textwidth}
    Comando PUMP2 [\unit{\percent}] & Contas escritas & FT2 [\unit{\litre\per\minute}] & LT [\unit{\percent}] & Altura $h$ [\unit{\milli\metre}] \\
    \num{25} & \vazio & & & \\
    \num{50} & \vazio & & & \\
    \num{75} & \vazio & & & \\
  \end{tabelaguia}
\end{table}

\begin{table}[H]
  \centering
  \caption{Calibração do transmissor de nível \texttt{LT} --- preencher durante a prática, com o tanque
    esvaziando em passos e a vazão nula em cada medição.}
  \label{tab:calibracao-lt}
  \begin{tabelaguia}{colspec={cccc|cccc},width=\textwidth}
    Ponto & $h$ [\unit{\milli\metre}] & Contas & $V_{LT}$ [\unit{\volt}] & Ponto & $h$ [\unit{\milli\metre}] & Contas & $V_{LT}$ [\unit{\volt}] \\
    1  & \vazio & & & 11 & & & \\
    2  & \vazio & & & 12 & & & \\
    3  & \vazio & & & 13 & & & \\
    4  & \vazio & & & 14 & & & \\
    5  & \vazio & & & 15 & & & \\
    6  & \vazio & & & 16 & & & \\
    7  & \vazio & & & 17 & & & \\
    8  & \vazio & & & 18 & & & \\
    9  & \vazio & & & 19 & & & \\
    10 & \vazio & & & 20 & & & \\
  \end{tabelaguia}
\end{table}

\section{Códigos de Exemplo}
Esta seção trás três códigos com exemplos da operação do CLP a partir de códigos Python.

\codigoarquivo[language=Python,
  caption={Conversões entre contas do CLP e grandezas físicas (arquivo
  \texttt{aulas/aula-01/src/conversoes.py}).},
  label={cod:conversoes}]{aulas/aula-01/src/conversoes.py}

\codigoarquivo[language=Python,
  caption={Leitura periódica das \texttt{tags} em unidade de engenharia (arquivo
  \texttt{aulas/aula-01/src/ler\_planta.py}).},
  label={cod:leitura}]{aulas/aula-01/src/ler_planta.py}

\codigoarquivo[language=Python,
  caption={Comando da válvula e da bomba em porcentagem, com intertravamento (arquivo
  \texttt{aulas/aula-01/src/comanda\_planta.py}).},
  label={cod:comando}]{aulas/aula-01/src/comanda_planta.py}


\section{Pesquisa}

\begin{roteiro}
  \item Calcule a resolução do conversor A/D --- \qty{16}{\bits} na faixa de \qtyrange{-10,5}{+10,5}{\volt}
    --- e converta-a para a unidade de engenharia de cada transmissor da Tabela~\ref{tab:entradas}. Discuta se essa resolução é limitante para a medição de vazão do circuito frio.
  \item Mostre o que aconteceria, no CLP, se o programa escrevesse \num{32768} numa \sw{tag} de saída, e aponte a linha de \texttt{conversoes.py} que impede esse erro.

  \item A Tabela~\ref{tab:entradas} traz a conversão \textbf{nominal} de \cd{LT}, em porcentagem de enchimento. Explique por que ela \textbf{não substitui} a curva de calibração levantada nam Tabela~\ref{tab:calibracao-lt} e cite duas fontes de erro que a calibração revela e a conversão nominal esconde.
  
  \item O tanque de processo é cilíndrico. A partir do seu diâmetro interno e da reta de calibração obtida nam Tabela~\ref{tab:calibracao-lt}, escreva a relação entre a leitura de \cd{LT}, em volts, e o volume de líquido, em litros, e discuta que hipóteses esta conversão está assumindo.

\end{roteiro}

\begin{comment}

  
\section{Sugestão de Redação do Relatório}

Cada aula corresponde a uma seção do artigo entregue ao final da disciplina. Para esta aula, a seção é
predominantemente descritiva: o objetivo é apresentar a planta e demonstrar que a cadeia de comunicação está
íntegra nos dois sentidos, não ainda extrair modelos ou sintonizar controladores.

\subtituloguia{Introdução da seção}

Apresente em poucas frases a planta TQ CE117, o CLP e a biblioteca \sw{pycomm3}, e o objetivo desta etapa:
estabelecer e validar experimentalmente o caminho de dados entre o processo e o programa em Python, sobre o
qual as demais etapas do trabalho serão construídas.

\subtituloguia{Metodologia}

Descreva, em texto corrido e na terceira pessoa (não copie o roteiro em imperativo), como a sessão
\sw{EtherNet/IP} com o CLP foi aberta, quais \sw{tags} foram lidas e escritas, como as conversões A/D e D/A
foram implementadas em \texttt{conversoes.py} e como o ensaio de enchimento foi conduzido. Inclua um trecho
comentado do código de aquisição e explique o papel de cada função. Registre explicitamente o intertravamento
adotado entre a válvula \cd{S} e \cd{PUMP2} e onde ele está implementado. Descreva ainda o procedimento de
calibração do transmissor de nível: como a altura foi medida, por que cada ponto foi tomado com vazão nula e
com a válvula de dreno fechada, e quantos pontos cobrem a faixa ensaiada.

\subtituloguia{Resultados}

\begin{itemize}
  \item Reproduza a Tabela~\ref{tab:leituras} acrescentando uma coluna com a tensão calculada à mão pela
    \req{eq:adc} e o desvio percentual em relação ao valor apresentado pelo programa. Esse desvio é o
    indicador de que a conversão foi implementada corretamente.
  \item A partir da Tabela~\ref{tab:enchimento}, construa um gráfico com o comando de \cd{PUMP2}, em
    porcentagem, no eixo horizontal e a vazão \cd{FT2}, em \unit{\litre\per\minute}, no eixo vertical. Indique se a relação
    observada é monotônica e se há zona morta em comandos baixos.
  \item Construa o gráfico de calibração de \cd{LT} a partir da Tabela~\ref{tab:calibracao-lt}, com a altura
    $h$, em \unit{\milli\metre}, no eixo horizontal e a tensão lida, em \unit{\volt}, no eixo vertical.
    Marque os \num{20} pontos medidos e ajuste sobre eles uma reta $V_{LT} = a\,h + b$ por mínimos quadrados.
  \item Informe os coeficientes obtidos com suas unidades --- a sensibilidade $a$ em \unit{\volt\per\milli\metre}
    e o \sw{offset} $b$ em \unit{\volt} ---, o coeficiente de determinação $R^2$ do ajuste e o maior desvio
    entre um ponto medido e a reta. Apresente também a relação inversa $h = (V_{LT} - b)/a$, que é a forma
    efetivamente usada pelo programa para converter a leitura em altura.
\end{itemize}

\subtituloguia{Discussão}

Comente a integridade da cadeia de medição: se os valores convertidos são compatíveis com o estado físico
observado da planta (tanque vazio, bomba parada) e a que fatores atribuir eventuais discrepâncias --- erro no
fator de escala, ruído do transmissor, resolução do conversor, atraso entre requisições CIP. Discuta a qualidade da calibração de \cd{LT}: se os resíduos em torno da reta ajustada se distribuem ao
acaso ou seguem um padrão --- uma curvatura sistemática apontaria não linearidade do transmissor, enquanto
uma dispersão uniforme aponta a incerteza da leitura da régua ---, e estime a resolução em altura que a
sensibilidade obtida permite, em \unit{\milli\metre} por conta. Comente também o
comportamento do circuito frio sob comandos crescentes de \cd{PUMP2}, apontando indícios de zona morta,
saturação ou atraso de transporte. Encerre discutindo o período de amostragem que o acesso via \sw{pycomm3}
permite sustentar e o que isso significa para o laço de controle a ser fechado nas aulas seguintes.
\end{comment}



%%%%%%%%%%%%%%%%%%%%%%%%%%%%%%%%%%%%%%%%%%%%%%%%%%%%%%%%%%%%%%%%%%%%%%%%%%%%%%%%
%% Aula exemplo -- vitrine dos recursos da classe guiapratico
%%
%% Este arquivo nao pretende ser uma aula de verdade: ele exercita, um a um, os
%% comandos e ambientes da classe, para servir de consulta rapida na hora de
%% escrever as aulas reais. Copie daqui os trechos que interessarem.
%%
%% ARMADILHA UNICA: \cd e verbatim (\lstinline). Ele NAO pode aparecer dentro
%% do argumento de outra macro -- \objetivos, \campo, \obs, \caption, \section,
%% \item de \objetivos etc. Nesses lugares use \texttt{\textbackslash comando}.
%%%%%%%%%%%%%%%%%%%%%%%%%%%%%%%%%%%%%%%%%%%%%%%%%%%%%%%%%%%%%%%%%%%%%%%%%%%%%%%

% \aula[titulo curto para o sumario]{Titulo completo}
% Abre pagina nova, imprime a faixa cinza "Aula N" com o brasao e zera os
% contadores de secao, figura, tabela, equacao e codigo.
\aula[Recursos do template]{Recursos do Template --- Aula de Exemplo}

% ---------------------------------------------------------------- cabecalho --
% Campos rotulados do topo da aula. Um campo vazio simplesmente nao e impresso.

\objetivos{%
  \item conhecer os comandos de estrutura da classe
        (\texttt{\textbackslash aula}, \texttt{\textbackslash section},
        \texttt{\textbackslash secaoguia}, \texttt{\textbackslash subtituloguia});
  \item exercitar os blocos próprios do guia --- roteiro, tabela de
        preenchimento, figura, listagem de código e citação;
  \item servir de referência rápida durante a escrita das demais aulas.}

\listamaterial{1 computador com \sw{MATLAB}/\sw{Octave} ou \sw{Python};
  1 osciloscópio digital; 1 gerador de funções; 1 protoboard; 1 circuito
  integrado \ci{SN7400}; resistores e capacitores diversos; cabos banana--jacaré.}

% \campo{Rotulo:}{conteudo} -- campo rotulado generico, para quando os campos
% prontos (\objetivos, \listamaterial) nao bastarem.
\campo{Duração:}{2 horas/aula (1h40min).}

% ------------------------------------------------------------------ texto ----
% \introducao imprime o rotulo "Introducao:" (e um \subtituloguia pronto).
\section{Introdução}

Este documento reúne, em um só lugar, tudo o que a classe \cd{guiapratico}
oferece para a escrita de uma aula prática. O texto corrido é composto em
DejaVu Sans, com parágrafos justificados; termos estrangeiros e nomes de
programas passam por \cd|\sw{...}| --- como em \sw{software}, \sw{datasheet} e
\sw{Simulink} --- e códigos de circuitos integrados por \cd|\ci{...}|, que
impede a quebra de linha no meio da sigla: \ci{SN7400}, \ci{CD4017}.

Trechos curtos de código no meio da frase saem com \cd|\cd{...}|: chame
\cd{step(sys)} para obter a resposta ao degrau. Se o trecho contiver \verb|%| ou
chaves desbalanceadas, troque as chaves por outro delimitador, como em
\verb+\cd|a % b|+.

A referência teórica da disciplina é \textcite{ogata2010}; para o projeto de
controladores no domínio da frequência, veja também \cite{nise2017,dorf2018}.
Uma citação com página fica assim: \cite[p.~42]{franklin2019}. Citação de
citação usa \texttt{\textbackslash apud}: \apud{astrom2004}{ogata2010}.

% ATENCAO: \cd nao funciona dentro do argumento de \obs -- use \texttt.
\obs{o comando \texttt{\textbackslash obs} produz este parágrafo iniciado por
\textbf{OBS --}, útil para avisos de segurança ou pegadinhas da montagem.}

% -------------------------------------------------------------- estrutura ----
% Duas familias de titulos convivem no mesmo texto:
%   \section/\subsection/\subsubsection -> numerados (1.1, 1.1.1, 1.1.1.1)
%   \secaoguia/\subtituloguia           -> sem numero, fora do sumario

\section{Títulos numerados}\label{sec:titulos}

Esta é uma \cd{\section}: numerada com o número da aula como prefixo
(\ref{sec:titulos}) e referenciável por \cd|\label|/\cd|\ref|. Por padrão as
seções não aparecem no sumário --- para incluí-las, aumente o \cd{tocdepth} no
preâmbulo do \cd{main.tex}.

\subsection{Subseção}

Um nível abaixo. A numeração acompanha: a Seção~\ref{sec:titulos} gera esta
subseção \thesubsection.

\subsubsection{Sub-subseção}

O nível mais profundo numerado pela classe.

\secaoguia{Título de bloco sem número}

O \cd{\secaoguia} imprime um título em caixa alta, sem numeração e fora do
sumário --- é o mesmo estilo usado em \textbf{REFERÊNCIAS} e nos apêndices.

\subtituloguia{Subtítulo em negrito}

Já o \cd{\subtituloguia} produz um subtítulo indentado, também sem número.

% ---------------------------------------------------------------- formulas ---
\section{Equações}

O \cd{amsmath} já vem carregado pela classe. O contador de equações é zerado a
cada aula (as figuras, tabelas e listagens é que levam o número da aula como
prefixo), e a equação aceita rótulo:

\begin{equation}\label{eq:pid}
  u(t) = K_p\,e(t) + K_i \int_{0}^{t} e(\tau)\,\mathrm{d}\tau
       + K_d \frac{\mathrm{d}e(t)}{\mathrm{d}t}.
\end{equation}

Em Laplace, a \req{eq:pid} vira a função de transferência

\begin{equation}\label{eq:pid-s}
  C(s) = K_p + \frac{K_i}{s} + K_d s
       = K_p\left(1 + \frac{1}{T_i s} + T_d s\right),
\end{equation}

com $T_i = K_p/K_i$ e $T_d = K_d/K_p$. Matemática no meio do texto também
funciona: o sistema de primeira ordem $G(s) = K/(\tau s + 1)$ tem constante de
tempo $\tau$ e ganho estático $K$.

Para citar uma equação no texto use \cd|\req{rotulo}|, que escreve
\enquote{Equação} seguida do número entre parênteses, como na \req{eq:pid-s}.

% ----------------------------------------------------------------- unidades ---
\section{Números e unidades}

A classe carrega o \cd{siunitx}, que é a forma padrão de escrever grandezas
neste guia --- nunca digite a unidade à mão. Ele já vem configurado para o
português: vírgula decimal, milhar separado por espaço fino e faixas escritas
com \enquote{a}.

\begin{itemize}
  \item \cd|\qty{500}{\milli\second}| $\rightarrow$ \qty{500}{\milli\second}
        --- valor com unidade (o espaço correto entre eles é automático e
        não quebra linha);
  \item \cd|\qtyrange{-10,5}{+10,5}{\volt}| $\rightarrow$
        \qtyrange{-10,5}{+10,5}{\volt} --- faixa com a unidade escrita uma
        única vez;
  \item \cd|\unit{\litre\per\minute}| $\rightarrow$ \unit{\litre\per\minute}
        --- só a unidade, para cabeçalhos de tabela e rótulos de eixo;
  \item \cd|\num{-32768}| $\rightarrow$ \num{-32768} --- só o número, com o
        sinal e o agrupamento de milhar corretos;
  \item \cd|\qty{28,3}{\percent}| $\rightarrow$ \qty{28,3}{\percent} ---
        porcentagem (repare que a vírgula decimal é escrita como vírgula).
\end{itemize}

Todos funcionam dentro e fora do modo matemático. Unidades fora do SI usadas
no guia são declaradas na classe com \cd|\DeclareSIUnit|: \cd|\bit|/\cd|\bits|, \cd|\rpm|
e \cd|\voltpp| ($\qty{5}{\voltpp}$, por exemplo).

% ----------------------------------------------------------------- figuras ---
\section{Figuras}

\cd|\figuraguia[largura]{arquivo}{legenda}| coloca a figura centrada, com
legenda numerada, buscando o arquivo em \cd{img/} \emph{ao lado} deste
\cd{.tex} --- não importa de onde o arquivo seja incluído.

\figuraguia[0.8\textwidth]{exemplo.png}{Imagem de exemplo, incluída com
  \texttt{\textbackslash figuraguia}.}

Para uma imagem compartilhada entre várias aulas, guardada em \cd{figuras/} na
raiz do projeto, use \cd|\figuraglobal{figuras/arquivo.pdf}{legenda}|. E para
inserir a imagem sem \sw{float} nem legenda --- dentro de uma tabela, por
exemplo --- existe \cd|\includefigure[largura]{arquivo}|.

% Figura montada "na mao" em vez de \figuraguia, para que o \label venha logo
% depois do \caption -- so assim ela pode ser referenciada por \ref (o
% \figuraguia nao tem onde encaixar o rotulo).
\begin{figure}[H]
  \centering
  \includefigure[0.5\textwidth]{exemplo.png}
  \caption{Circuito usado na prática; aqui o \texttt{\textbackslash label} vem
    logo após o \texttt{\textbackslash caption}, o que permite a referência
    cruzada feita mais adiante, no roteiro.}
  \label{fig:exemplo-ref}
\end{figure}

A chamada no texto é feita com \cd|\rfig{rotulo}|, que escreve
\enquote{Figura} seguida do número --- por exemplo, \rfig{fig:exemplo-ref}.

% ----------------------------------------------------------------- tabelas ---
\section{Tabelas}

O ambiente \cd{tabelaguia} desenha a tabela no estilo do guia: fios pretos e
primeira linha em cinza, negrito. A especificação das colunas segue o
\cd{tabularray} (\cd{c}, \cd{l}, \cd{r}, \cd{X}\ldots).

\begin{table}[H]
  \centering
  \caption{Tabela verdade da porta NAND, preenchida pelo professor.}
  \label{tab:verdade}
  \begin{tabelaguia}{colspec={ccc},width=0.55\textwidth}
    $A$ & $B$ & $S = \overline{A \cdot B}$ \\
    0 & 0 & 1 \\
    0 & 1 & 1 \\
    1 & 0 & 1 \\
    1 & 1 & 0 \\
  \end{tabelaguia}
\end{table}

Nas tabelas de preenchimento à mão, \cd{\vazio} dá altura útil à célula vazia:

\begin{table}[H]
  \centering
  \caption{Medidas do ensaio --- preencher durante a prática.}
  \label{tab:medidas}
  \begin{tabelaguia}{colspec={cccc},width=0.85\textwidth}
    Ensaio & $V_{\text{ent}}$ [\unit{\volt}] & $V_{\text{sai}}$ [\unit{\volt}] & $t_s$ [\unit{\milli\second}] \\
    1 & \vazio & & \\
    2 & \vazio & & \\
    3 & \vazio & & \\
  \end{tabelaguia}
\end{table}

A Tabela~\ref{tab:medidas} também pode ser citada por \cd{\ref}, como a
Tabela~\ref{tab:verdade}.

% ---------------------------------------------------------------- roteiros ---
% \pratica imprime o titulo "Pratica:" alinhado a esquerda.
\pratica

O ambiente \cd{roteiro} numera os passos no formato \textbf{1}~-- e aceita um
nível de subitens (\textbf{1.1}~--).

\begin{roteiro}
  \item Monte na protoboard o circuito da \rfig{fig:exemplo-ref}, com o
        \ci{SN7400} alimentado em \qty{5}{\volt}.

  \item Ajuste o gerador de funções para onda quadrada de \qty{1}{\kilo\hertz},
        amplitude \qty{5}{\voltpp} e nível DC de \qty{2,5}{\volt}:
        \begin{roteiro}
          \item conecte a saída do gerador à entrada $A$;
          \item mantenha a entrada $B$ em nível lógico alto;
          \item verifique a forma de onda no osciloscópio antes de prosseguir.
        \end{roteiro}

  \item Registre na Tabela~\ref{tab:medidas} os valores medidos.

  \item Compare o resultado com a Tabela~\ref{tab:verdade} e justifique as
        divergências, se houver.
\end{roteiro}

% ---------------------------------------------------------------- listagens --
\section{Listagens de código}

Blocos de código têm realce de sintaxe, numeração de linhas e legenda
\cd{Codigo N.M:} --- numerada por aula, como as figuras. Acentuação dentro das
listagens funciona normalmente.

\begin{codigomatlab}[caption={Resposta ao degrau de uma planta de segunda ordem.},
                     label={cod:degrau}]
% Planta de segunda ordem subamortecida
wn   = 2*pi*5;        % frequência natural [rad/s]
zeta = 0.3;           % coeficiente de amortecimento
G    = tf(wn^2, [1 2*zeta*wn wn^2]);

% Malha fechada com controlador proporcional
Kp = 4;
T  = feedback(Kp*G, 1);

step(T); grid on;
title('Resposta ao degrau em malha fechada');
\end{codigomatlab}

O Código~\ref{cod:degrau} usa o ambiente \cd{codigomatlab}. Há ainda
\cd{codigopython}, \cd{codigoc}, \cd{codigocpp}, \cd{codigoarduino},
\cd{codigovhdl}, \cd{codigobash} e \cd{codigotexto} (sem realce). O ambiente
genérico \cd{codigo} aceita \cd{language=} no argumento opcional, junto de
qualquer outra chave do \cd{listings}.

\begin{codigoarduino}[caption={Leitura do sensor e acionamento do atuador.},
                      label={cod:arduino}]
const int PINO_SENSOR  = A0;
const int PINO_ATUADOR = 9;

void setup() {
  Serial.begin(115200);
  pinMode(PINO_ATUADOR, OUTPUT);
}

void loop() {
  int leitura = analogRead(PINO_SENSOR);   // leitura de 0 a 1023
  analogWrite(PINO_ATUADOR, leitura / 4);  // saída PWM de 0 a 255
  Serial.println(leitura);
  delay(10);
}
\end{codigoarduino}

Sem numeração de linhas e sem realce, para saída de terminal:

\begin{codigotexto}[caption={Saída esperada no monitor serial.}, style=semnumeros]
$ python degrau.py
K = 2.00, tau = 3.00 s, theta = 1.50 s
\end{codigotexto}

\subtituloguia{Código em arquivo externo}

Melhor que copiar e colar: o fonte fica em \cd{src/}, roda de verdade e o guia
mostra sempre a versão atual. As chaves \cd{firstline}/\cd{lastline} recortam
um trecho, e \cd{style=semnumeros} desliga a numeração de linhas.

\codigoarquivo[language=Python, firstline=14, lastline=21,
  caption={Identificação pelo método dos dois pontos (arquivo
  \texttt{aulas/aula-exemplo/src/degrau.py}).},
  label={cod:identifica}]{aulas/aula-exemplo/src/degrau.py}

O Código~\ref{cod:identifica} implementa o método de Smith, discutido em
\textcite{astrom2008pid}.

% ------------------------------------------------------------------ listas ---
\section{Listas de apoio}

Além do \cd{roteiro}, valem as listas padrão do \LaTeX, já compactadas pela
classe:

\begin{itemize}
  \item itens sem numeração, para enumerar materiais ou cuidados;
  \item aceitam aninhamento:
    \begin{itemize}
      \item segundo nível;
    \end{itemize}
\end{itemize}

\begin{enumerate}[label=\alph*),leftmargin=1.2cm]
  \item o \cd{enumitem} está carregado, então rótulos personalizados funcionam;
  \item útil para as alíneas exigidas pela ABNT no relatório.
\end{enumerate}

% ---------------------------------------------------------------- pesquisa ---
% \pesquisa imprime o titulo "PESQUISA" centrado.
\pesquisa

Bloco para a parte teórica que o aluno responde em casa.

\begin{roteiro}
  \item Consulte o \sw{datasheet} do \ci{SN7400} \cite{sn7400} e informe a
        corrente máxima de saída em nível baixo ($I_{OL}$).

  \item Explique, com suas palavras, por que o método dos dois pontos usa
        justamente \qty{28,3}{\percent} e \qty{63,2}{\percent} do valor final.

  \item Pesquise na documentação do \sw{Control System Toolbox}
        \cite{mathworks-cst} qual função retorna as margens de ganho e de fase.
\end{roteiro}

\obs{a fonte de toda pesquisa deve ser citada no relatório, conforme a
  \cite{nbr6023} --- veja o Apêndice~B.}

% --- referências (NBR 14724)
\referencias

% --- apêndices
%%%%%%%%%%%%%%%%%%%%%%%%%%%%%%%%%%%%%%%%%%%%%%%%%%%%%%%%%%%%%%%%%%%%%%%%%%%%%%%%
%% Apêndice A -- Normas de funcionamento do laboratório
%%%%%%%%%%%%%%%%%%%%%%%%%%%%%%%%%%%%%%%%%%%%%%%%%%%%%%%%%%%%%%%%%%%%%%%%%%%%%%%
\apendice[Normas de Funcionamento do Laboratório]%
        {Normas de Funcionamento do Laboratório\\ de Sistemas Digitais e Microprocessadores}

A supervisão dos laboratórios, visando à aprendizagem e por questões de segurança,
estabelece para as aulas práticas e demais atividades no Laboratório de Eletrônica
Analógica e Digital um conjunto de Normas para Aulas Práticas e de Segurança, listadas
a seguir.

\secaoguia{I -- Normas gerais}

\begin{normas}[I]
  \norma{OBJETOS} --- os objetos volumosos de uso pessoal (mochilas, pastas, livros,
  cadernos etc.) devem ser colocados nos escaninhos disponíveis, na frente da sala ou
  sob as bancadas do laboratório, mas nunca sobre as mesmas;

  \norma{CELULAR} --- NÃO é permitido o uso do telefone celular ligado durante as aulas
  práticas. Recomenda-se desligá-lo antes de cada aula;

  \norma{PONTUALIDADE} --- para as aulas de laboratório existe uma tolerância de até 15
  minutos. \textbf{O aluno que chegar após este período não poderá assistir às aulas e
  só poderá efetuar a sua reposição por motivo justificável (por escrito).}

  \norma{NÚMERO DE ALUNOS POR BANCADA} --- as montagens dos circuitos serão feitas por
  até 4 (quatro) alunos por bancada, divididos em grupos de dois alunos, que deverão
  fazer o relatório da aula por grupo e entregá-lo na semana seguinte à aula;

  \norma{POSTURA PROFISSIONAL:}
  \begin{enumerate}[label=\alph*),leftmargin=1.2cm]
    \item não será permitido aos alunos assistirem aula trajando bermuda e/ou
          camisetas, por medidas de segurança e por aspectos de postura profissional;
    \item o laboratório não é um lugar de bate-papo ou de perda de tempo;
    \item não é permitido ao aluno fumar ou fazer lanche no laboratório;
  \end{enumerate}
\end{normas}

\secaoguia{II -- Normas de organização e segurança}

\begin{normas}[II]
  \norma{PROCEDIMENTO INICIAL} --- antes de cada montagem, ligue o disjuntor
  diferencial da bancada. Após o término da mesma, desligue-o;

  \norma{DISPOSIÇÃO DOS COMPONENTES E EQUIPAMENTOS} --- ao término de cada montagem, os
  equipamentos e componentes devem ser dispostos como recebidos antes do início da aula;

  \norma{PARA QUALQUER PROCEDIMENTO} --- execute-o somente se estiver habilitado. Na
  DÚVIDA, informe-se com o professor ou com o laboratorista;

  \norma{SOLICITE A PRESENÇA DO PROFESSOR} em sua bancada, no caso de dúvidas ou de
  acidentes.
\end{normas}


%%%%%%%%%%%%%%%%%%%%%%%%%%%%%%%%%%%%%%%%%%%%%%%%%%%%%%%%%%%%%%%%%%%%%%%%%%%%%%%%
%% Apêndice B -- Roteiro para elaboração dos relatórios
%%%%%%%%%%%%%%%%%%%%%%%%%%%%%%%%%%%%%%%%%%%%%%%%%%%%%%%%%%%%%%%%%%%%%%%%%%%%%%%
\apendice{Roteiro para Elaboração dos Relatórios}

\subtituloguia{1) Instruções gerais}

O relatório deverá ser feito à mão. Erros no uso correto do Português são penalizados
com redução de nota. A parte inicial do relatório deverá conter:

\begin{enumerate}[label=\alph*),leftmargin=1.2cm]
  \item nome e número de matrícula dos alunos que realizaram a prática;
  \item o número e o título da prática;
  \item a data em que se realizou a prática.
\end{enumerate}

\subtituloguia{2) Introdução teórica}

Esta seção deverá conter uma introdução sobre a prática realizada, com enfoque no
objetivo e no conteúdo que foi avaliado experimentalmente. Expresse os principais
circuitos relacionados à prática, bem como os dados obtidos por simulação e/ou
experimentalmente e os calculados. Liste todos os materiais utilizados na prática.

\subtituloguia{3) Memória de cálculo}

Nesta seção deverão ser apresentadas as questões, tabelas e/ou gráficos relativos à
parte Prática de cada um dos experimentos. É permitido responder no próprio roteiro da
aula prática e anexá-lo ao relatório.

\subtituloguia{4) Resultado da pesquisa}

Nesta seção deverá ser apresentado o assunto pesquisado, caso exista esse tópico no
roteiro da prática realizada. Não se esqueça de colocar a fonte da pesquisa (livro,
internet, artigos).

\subtituloguia{5) Conclusão}

Nesta seção deve-se elaborar e redigir as principais conclusões observadas pelo grupo.
A conclusão deve versar sobre os objetivos. É importante ressaltar que a conclusão é a
parte mais importante do relatório.

%\apendice[Mapa da Teoria de Controle]{Mapa da Teoria de Controle}

A \rfig{fig:f9k3d8} ilustra um "mapa" da Teoria de Controle, indicando seus diferentes campos e as principais técnicas e ferramentas que as compõe. 
%


\begin{figure}[H]
    \centering
    \includefigure[\textwidth]{mapa-controle.png}
    \caption{Mapa de controle}
    \label{fig:f9k3d8}
\end{figure}


\end{document}
